\begin{quadro*}[tbp]
\caption{Condições necessárias para interpretar e utilizar comparações entre métodos}\label{qua:aplicacao}
\begingroup\singlespacing\fontsize{10}{12}\selectfont
\begin{tabularx}{\textwidth}{|P{.23\textwidth}|P{.34\textwidth}|X|}
\hline\textbf{Condicionante} & \textbf{Informação a preservar} & \textbf{Implicação para uso dos dados}\\\hline
Aquecimento e semivoláteis & Temperaturas, secagem e retenção ou compensação de material nos dois sistemas \parencite{rsl201,rsl101,rsl144}. & Interpretar a diferença em relação ao tratamento da amostra, incluindo os artefatos do filtro.\\\hline
Água e propriedades ópticas & Umidade, composição, configuração e identificação das séries bruta e corrigida \parencite{rsl170,rsl089,rsl388}. & Delimitar em quais condições uma correção de umidade melhora a relação, sem assumir resposta única por tecnologia.\\\hline
Seleção por tamanho e conversão em massa & Fração, entrada, densidade e forma assumidas na conversão \parencite{rsl136,rsl288,rsl283}. & Evitar transportar desempenho de uma fração ou comparador para outro sem avaliação.\\\hline
Concentração e definição da estatística & Diferenças absolutas e relativas, intercepto, dispersão e tipo de razão \parencite{rsl073,rsl129,rsl085}. & Examinar baixas concentrações e distinguir associação de concordância.\\\hline
Pareamento e escala temporal & Intervalos efetivamente coincidentes, cobertura, exclusões e denominador por contraste \parencite{rsl059,rsl026}. & Separar comportamento no intervalo de coleta e resultado agregado; registrar rupturas instrumentais na série.\\\hline
Desenvolvimento e transferência de correções & Dados de desenvolvimento, dados de avaliação, forma do modelo e domínio ambiental \parencite{rsl029,rsl085,rsl219,rsl342}. & Restringir a aplicação ao alcance demonstrado e examinar o desempenho em condições distintas do ajuste.\\\hline
\end{tabularx}
\Fonte{Elaboração própria a partir dos condicionantes e desenhos discutidos nos relatos citados.}
\NotaQuadro{O quadro organiza a interpretação e não estabelece um protocolo novo de equivalência ou uma classificação de risco dos estudos.}
\endgroup
\end{quadro*}
